\section{Factor-ray setup and source ledgers}
\label{sec:setup}

Fix numbers $0<\alpha<\beta$ such that
$\supp W\subset[\alpha,\beta]$.  The shift is an integer, and $X$ is
always large enough that every target selected by the weights is positive.
All sums are finite.

\subsection{Ray compression and the inherited spectral identities}

For every factor pair, write uniquely
\begin{equation}
 n=ak,\qquad m=bk,\qquad (a,b)=1,\qquad k=(m,n).
 \label{eq:ray-coordinates}
\end{equation}
Then $mn=abk^2$ and $n/m=a/b$, so substitution gives
\eqref{eq:intro-ray-compression}--\eqref{eq:intro-ray-coefficient}.
The raw ray coefficient is
\begin{equation}
 c^{\mathrm{raw}}_{a,b}(h;X,D)=
 \sum_{\substack{k\geq1\\ak>D}}
 u(ak)\Lambda(abk^2+h)W(abk^2/X),
 \label{eq:raw-ray-coefficient}
\end{equation}
and its squared energy is denoted by $\cE^{\mathrm{raw}}_{h,D}(X)$.

Two geometric facts from TPC-11 will be used below.  We record them to
fix all normalizations.  Distinct primitive rays in the support satisfy
\begin{equation}
 \left|\log\frac ab-\log\frac cd\right|
 \geq
 \delta_X,
 \qquad
 \delta_X=2\arsinh\!\left(\frac1{2\beta X}\right).
 \label{eq:ray-spacing}
\end{equation}
Indeed, $|ad-bc|\geq1$ and $abcd\leq(\beta X)^2$; the identity
\[
 2\sinh\!\left(\frac12\left|\log\frac{ad}{bc}\right|\right)
 =\frac{|ad-bc|}{\sqrt{abcd}}
\]
proves the claim.

The generalized Hilbert inequality of \citet{MontgomeryVaughan1974}
therefore gives, for every bounded interval $I\subset\R$,
\begin{equation}
 \int_I|\cR_{h,D}(X;t)|^2dt
 =|I|\cE_{h,D}(X)+O_W\bigl(X\cE_{h,D}(X)\bigr).
 \label{eq:MV-ray-law}
\end{equation}
The same identity holds for the raw transform.  With the Fourier convention
$\widehat f(\xi)=\int_\R f(t)e^{-it\xi}\,dt$ and
\begin{equation}
 \kappa(t)=\frac2\pi\left(\frac{\sin(t/2)}t\right)^2,
 \qquad
 \widehat\kappa(\xi)=(1-|\xi|)_+,
 \label{eq:Fejer-kernel}
\end{equation}
Fourier expansion gives the exact formula
\begin{equation}
 \int_\R\kappa(t/T)|\cR_{h,D}(X;t)|^2dt
 =T\cE_{h,D}(X)
 \quad\text{whenever }T\delta_X\geq1.
 \label{eq:exact-Fejer-law}
\end{equation}
The endpoint value of $\widehat\kappa$ is zero, so equality is valid also
when a scaled gap equals one.

\subsection{Weights on product centers}

For an integer $r\geq1$, define
\begin{align}
 A_D(r)
 &:={}|W(r/X)|^2
 \sum_{\substack{n\mid r\\n>D}}u(n)^2,
 \label{eq:product-diagonal-weight}\\
 B_{D,t}(r)
 &:={}W(r/X)
 \sum_{\substack{n\mid r\\n>D}}
 u(n)(n^2/r)^{it}.
 \label{eq:product-signed-weight}
\end{align}
The first is the source weight on the target-square diagonal.  The second
is the coefficient that appears after the first-moment transform is
grouped by $r=mn$.

\begin{lemma}[Pointwise and quadratic product ledgers]
\label{lem:product-ledgers}
Uniformly in $D\geq1$ and $t\in\R$,
\begin{align}
 A_D(r)&\ll_W\one_{[\alpha X,\beta X]}(r)
 \bigl((\log(2r))^2+\tau(r)\bigr),
 \label{eq:A-pointwise}\\
 |B_{D,t}(r)|&\ll_W\one_{[\alpha X,\beta X]}(r)
 \bigl(\log(2r)+\tau(r)\bigr),
 \label{eq:B-pointwise}
\end{align}
and
\begin{align}
 \sum_{r\geq1}A_D(r)^2&\ll_WX(\log(2X))^4,
 \label{eq:A-second-moment}\\
 \sum_{r\geq1}|B_{D,t}(r)|^2&\ll_WX(\log(2X))^3.
 \label{eq:B-second-moment}
\end{align}
\end{lemma}

\begin{proof}
Since $(\Lambda-1)^2\leq2(1+\Lambda^2)$,
\[
 \sum_{d\mid r}u(d)^2
 \leq2\tau(r)+2\sum_{p^j\mid r}(\log p)^2.
\]
If $r=\prod_pp^{v_p}$, then
\[
 \sum_{p^j\mid r}(\log p)^2
 =\sum_pv_p(\log p)^2
 \leq\left(\sum_pv_p\log p\right)^2=(\log r)^2.
\]
This proves \eqref{eq:A-pointwise}.  Also
\[
 \sum_{d\mid r}|u(d)|
 \leq\sum_{d\mid r}\Lambda(d)+\tau(r)
 =\log r+\tau(r),
\]
which proves \eqref{eq:B-pointwise}, uniformly because the phase has
modulus one.  Finally use the classical estimate
$\sum_{r\leq y}\tau(r)^2\ll y(\log(2y))^3$ and the fixed support of $W$.
\end{proof}

The total mass of $A_D$ has an asymptotic that will provide the main
coefficient.

\begin{lemma}[Source diagonal]
\label{lem:source-diagonal}
For every fixed $\eps>0$, uniformly for $1\leq D\leq X^{1-\eps}$,
\begin{equation}
 \begin{aligned}
 \cS_D(X)
 &:={}\sum_{r\geq1}A_D(r)
 =\sum_{\substack{n>D\\m\geq1}}
 u(n)^2|W(mn/X)|^2\\
 &=\frac12XI_2(W)
 \bigl((\log X)^2-(\log D)^2\bigr)
 +O_W(X\log(2X)).
 \end{aligned}
 \label{eq:source-diagonal}
\end{equation}
\end{lemma}

\begin{proof}
Euler summation gives
\[
 \sum_{m\geq1}|W(mn/X)|^2
 =\frac XnI_2(W)+O_W(1).
\]
The classical zero-free-region prime number theorem and partial summation
give \citep[Chapter~5]{IwaniecKowalski2004}
\begin{equation}
 \sum_{n\leq y}\frac{(\Lambda(n)-1)^2}{n}
 =\frac12(\log y)^2-\log y+O(1),
 \label{eq:u-square-harmonic}
\end{equation}
while $\sum_{n\leq y}u(n)^2\ll y\log(2y)$.  Insert the Euler estimate,
subtract the range $n\leq D$, and absorb the fixed support endpoints into
$O_W(X\log X)$.
\end{proof}

\subsection{The same-ray off-diagonal ledger}

For later use put
\begin{equation}
 \begin{aligned}
 \cL_D(X):={}&
 \sum_{(a,b)=1}
 \sum_{\substack{k\ne\ell\\ak,a\ell>D}}
 |u(ak)u(a\ell)|\\
 &\qquad\times
 |W(abk^2/X)W(ab\ell^2/X)|.
 \end{aligned}
 \label{eq:source-offdiagonal-ledger}
\end{equation}

\begin{lemma}[Source-ray square and off-diagonal ledgers]
\label{lem:source-ray-ledgers}
Uniformly for $Y\geq2$,
\begin{equation}
 \sum_{ak^2\leq Y}\frac{u(ak)^2}{ak}
 \ll(\log(2Y))^2,
 \label{eq:source-ray-square}
\end{equation}
and, for every $D\geq1$,
\begin{equation}
 \cL_D(X)\ll_WX(\log(2X))^2.
 \label{eq:source-offdiagonal-bound}
\end{equation}
\end{lemma}

\begin{proof}
For \eqref{eq:source-ray-square}, the contribution of the constant part
of $u^2\leq2(1+\Lambda^2)$ is
\[
 \sum_{k\leq\sqrt Y}\frac1k
 \sum_{a\leq Y/k^2}\frac1a
 \ll(\log(2Y))^2.
\]
If $ak=p^j$, then $k=p^\ell$ and $a=p^{j-\ell}$.  Dropping the extra
constraint $p^{j+\ell}\leq Y$ gives
\[
 \sum_{p^j\leq Y}(j+1)\frac{(\log p)^2}{p^j}
 \ll(\log(2Y))^2.
\]

For the second assertion, let $K_{a,b}$ be the radial indices selected by
the two weights and put $x_k=|u(ak)W(abk^2/X)|$.  Their fixed
multiplicative support implies $|K_{a,b}|\ll_Wk$ for every
$k\in K_{a,b}$.  Hence
\[
 \sum_{k\ne\ell}x_kx_\ell
 \leq |K_{a,b}|\sum_kx_k^2
 \ll_W\sum_{k\in K_{a,b}}k x_k^2.
\]
For fixed $a,k$, the number of supported $b$ is
$O_W(X/(ak^2))$; the additive one is absorbed by support nonemptiness.
Summing first over $b$ gives
\[
 \cL_D(X)
 \ll_WX\sum_{ak^2\ll_WX}\frac{u(ak)^2}{ak},
\]
and \eqref{eq:source-ray-square} completes the proof.
\end{proof}
